\def\sectiontag{From a rectangular array to a mathematical object}
\section{What matrices are}

\begin{frame}{A matrix can play several roles}
  \begin{columns}[T,onlytextwidth]
    \begin{column}{0.44\textwidth}
      \[
      A=\begin{pmatrix}
      2&1\\[-1mm]
      -1&3
      \end{pmatrix}
      \]
    \end{column}
    \begin{column}{0.54\textwidth}
      \vspace{0.2cm}
      \pill{softblue}{data}\quad
      \pill{softgreen}{coefficients}\quad
      \pill{softorange}{relations}

      \vspace{0.45cm}
      \pill{softblue}{transformation}\quad
      \pill{softgreen}{geometric action}
    \end{column}
  \end{columns}

  \vspace{0.75cm}
  \takeaway{The entries alone do not tell the whole story; the interpretation comes from the problem.}
\end{frame}

\begin{frame}{The first skill is reading}
  \begin{center}
  \begin{tikzpicture}[node distance=0.8cm]
    \tikzset{q/.style={rounded corners=5pt,draw=line,very thick,fill=paper,minimum width=2.55cm,minimum height=0.9cm,align=center,font=\bfseries}}
    \node[q] (s) {What is\\the shape?};
    \node[q,right=of s] (e) {What do\\entries mean?};
    \node[q,right=of e] (r) {How should rows\\and columns be read?};
    \draw[-{Latex[length=3mm]},thick,blue] (s)--(e);
    \draw[-{Latex[length=3mm]},thick,blue] (e)--(r);
  \end{tikzpicture}
  \end{center}

  \vspace{0.8cm}
  \concept{Calculation comes after orientation.}
\end{frame}

\begin{frame}{General notation}
  \begin{columns}[T,onlytextwidth]
    \begin{column}{0.58\textwidth}
      \[
      A=\begin{pmatrix}
      a_{11}&a_{12}&\cdots&a_{1n}\\
      a_{21}&a_{22}&\cdots&a_{2n}\\
      \vdots&\vdots&\ddots&\vdots\\
      a_{m1}&a_{m2}&\cdots&a_{mn}
      \end{pmatrix}
      \]
    \end{column}
    \begin{column}{0.38\textwidth}
      \begin{block}{Shape}
        \[
        A\in\R^{\mblue{m}\times\mred{n}}
        \]
        \mblue{$m$ rows}\\
        \mred{$n$ columns}
      \end{block}
      \begin{exampleblock}{Address}
        \[
        a_{\mblue{i}\mred{j}}
        \]
        row $i$, column $j$
      \end{exampleblock}
    \end{column}
  \end{columns}
\end{frame}

\begin{frame}{Read one matrix carefully}
  \begin{columns}[T,onlytextwidth]
    \begin{column}{0.52\textwidth}
      \[
      M=\begin{pmatrix}
      \cellcolor{softblue}4&\cellcolor{softblue}-2&\cellcolor{softblue}0\\
      7&1&\cellcolor{softgreen}5
      \end{pmatrix}
      \]
    \end{column}
    \begin{column}{0.44\textwidth}
      \begin{block}{Shape}
        $M\in\R^{2\times3}$
      \end{block}
      \begin{exampleblock}{Selected entries}
        $m_{11}=4$\\
        $m_{12}=-2$\\
        $m_{23}=5$
      \end{exampleblock}
    \end{column}
  \end{columns}

  \vspace{0.35cm}
  \mutedtext{Blue: first row. Green: entry in row 2, column 3.}
\end{frame}

\begin{frame}{Rows and columns are different objects}
  \begin{columns}[T,onlytextwidth]
    \begin{column}{0.48\textwidth}
      \begin{block}{Row matrix}
        \[
        r=\begin{pmatrix}4&-2&0\end{pmatrix}
        \]
        \centering $1\times3$
      \end{block}
    \end{column}
    \begin{column}{0.48\textwidth}
      \begin{exampleblock}{Column matrix}
        \[
        c=\begin{pmatrix}4\\-2\\0\end{pmatrix}
        \]
        \centering $3\times1$
      \end{exampleblock}
    \end{column}
  \end{columns}

  \vspace{0.5cm}
  \takeaway{Same numbers, different shape, different algebraic role.}
\end{frame}

\begin{frame}{A small gallery of shapes}
  \begin{columns}[T,onlytextwidth]
    \begin{column}{0.23\textwidth}
      \centering\pill{softblue}{row}\\[0.3cm]
      $\begin{pmatrix}a&b&c\end{pmatrix}$\\[0.2cm]
      \mutedtext{$1\times3$}
    \end{column}
    \begin{column}{0.23\textwidth}
      \centering\pill{softgreen}{column}\\[0.3cm]
      $\begin{pmatrix}a\\b\\c\end{pmatrix}$\\[0.2cm]
      \mutedtext{$3\times1$}
    \end{column}
    \begin{column}{0.23\textwidth}
      \centering\pill{softorange}{square}\\[0.3cm]
      $\begin{pmatrix}a&b\\c&d\end{pmatrix}$\\[0.2cm]
      \mutedtext{$2\times2$}
    \end{column}
    \begin{column}{0.23\textwidth}
      \centering\pill{softred}{rectangular}\\[0.3cm]
      $\begin{pmatrix}a&b&c\\d&e&f\end{pmatrix}$\\[0.2cm]
      \mutedtext{$2\times3$}
    \end{column}
  \end{columns}

  \vspace{0.8cm}
  \concept{Shape determines which operations are possible.}
\end{frame}

\begin{frame}{Checkpoint: before touching the entries}
  \begin{columns}[T,onlytextwidth]
    \begin{column}{0.62\textwidth}
      \Large
      \begin{enumerate}
        \item name the shape;
        \item locate entries by indices;
        \item separate rows from columns;
        \item identify the intended interpretation.
      \end{enumerate}
    \end{column}
    \begin{column}{0.34\textwidth}
      \begin{alertblock}{Do not start with}
        ``Which formula should I use?''
      \end{alertblock}
      \begin{exampleblock}{Start with}
        ``What object am I holding?''
      \end{exampleblock}
    \end{column}
  \end{columns}
\end{frame}

% =============================================================================
